\section{Phase 0 --- Contract decisions and repository foundation}

\textbf{Objective.} Remove architectural ambiguity before public APIs or native backends accumulate assumptions.

\textbf{Why first.} The architecture establishes the broad boundaries, but several type, ownership, and portability decisions remain open. Implementing against unresolved semantics would create expensive rework in every backend.

\deliverables
\begin{enumerate}
  \item A decision register resolving or explicitly deferring: image identity, layer identity scope, image snapshot semantics, concurrent replacement, resource ownership, async completion, operation versioning, cross-language ABI, and supported initial platforms.
  \item A package dependency diagram with import rules and an automated dependency-cycle check.
  \item A repository layout for protocols, runtime, providers, operation adapters, samples, tests, benchmarks, and documentation.
  \item Command-line build/test entry points and CI jobs for the platforms that can be supported immediately.
  \item A contribution rule: no new provider or operation is registered without a contract and conformance tests.
\end{enumerate}

\exitcriteria
\begin{itemize}
  \item No unresolved decision blocks the first API implementation.
  \item The package graph has no circular dependency.
  \item CI builds the empty/minimal package skeleton and runs a trivial test on the declared baseline platform.
  \item The phase-0 decision register identifies assumptions and deferred questions without presenting them as settled facts.
\end{itemize}

\textbf{Do not do yet:} write GPU kernels, add a large algorithm catalog, or define a general-purpose plugin loader.

\phase{1 --- ImageProcessingProtocol domain contract}

\textbf{Objective.} Specify the primary library first: the image-processing vocabulary and contracts that all operation implementations must obey.

\textbf{Scope.} This phase produces the domain contract. It does not implement native compute providers and does not let generic execution types dictate image semantics.

\deliverables
\begin{enumerate}
  \item \textbf{Image model:} logical image identity, revision, canvas/spatial descriptor, ordered layer references, metadata, and snapshot behavior.
  \item \textbf{Layer model:} typed content categories for raster, mask, label map, and annotations; explicit layer properties and target selection.
  \item \textbf{Pixel semantics:} descriptor fields for dimensions, components, precision, channels, strides, numerical range, color/transfer function, alpha interpretation, and scientific/non-color data.
  \item \textbf{Operation object:} stable identifier, semantic version, immutable parameters, validation rules, input/output contract, numerical contract, alias/mutation declaration, and effect classification.
  \item \textbf{Operation implementation contract:} declares supported operation identities/versions, input/output constraints, precision guarantees, access patterns, and required generic execution capabilities.
  \item \textbf{Qualifying function contract:} defines structured eligible/rejected/preparation-required results and machine-readable reasons. The qualifier is owned by the image-processing domain contract.
  \item \textbf{Composition contract:} defines typed sequential composition and the behavior of a composite operation. Multi-input operations use explicit typed input records; no public graph engine is introduced.
  \item \textbf{Image result contract:} distinguishes producing a new logical image from atomically replacing an existing image's state. Neither promises physical in-place reuse.
\end{enumerate}

\exitcriteria
\begin{itemize}
  \item Every public concept has one definition and one owner.
  \item At least three worked operation contracts (for example, Gaussian blur, resize, and alpha-aware blend) demonstrate that numerical and layer semantics can be expressed without provider-specific names.
  \item A qualification example proves that the same operation can be accepted by one implementation and rejected by another with a precise reason.
  \item API usage sketches type-check conceptually and expose no Metal/CUDA/OpenCV types.
  \item The protocol has no dependency on application UI or a specific algorithm library.
\end{itemize}

\textbf{Gate.} Do not implement provider adapters until the operation and implementation contracts can express real differences in precision, format, border mode, resource access, and output guarantees.

\phase{2 --- GeneralExecutionProtocol generic execution contract}

\textbf{Objective.} Implement the lower-level execution dependency required by \IPP{}, while keeping it independent of image-processing terminology.

\textbf{Dependency constraint.} \GEP{} MUST NOT import \IPP{}, define Image/Layer/Pixel semantics, or know operation names such as blur, blend, resize, or segmentation.

\deliverables
\begin{enumerate}
  \item Provider identity/version and device discovery.
  \item Device and execution-context abstraction with capability reporting.
  \item Generic resource handles with ownership, access, lifetime, import/export, and completion dependencies.
  \item A minimal submission/completion contract that can represent asynchronous work without forcing Metal command buffers, CUDA streams, OpenCL queues, and D3D command lists into one false native abstraction.
  \item Optional capability interfaces for buffers, images/textures, external-memory interoperability, peer transfer, compilation, and specialized synchronization.
  \item Structured provider errors, device-loss behavior, cancellation rules, and diagnostics.
  \item A fake/test provider that exercises resource and submission lifecycle without claiming to be a real GPU implementation.
\end{enumerate}

\exitcriteria
\begin{itemize}
  \item \GEP{} builds and tests without referencing any image-processing symbol.
  \item A fake provider completes, fails, cancels, and releases resources according to the contract.
  \item Tests prove that resources are not released before submitted work completes.
  \item Native provider differences remain accessible behind provider-specific capabilities; no universal command-buffer fiction is introduced.
  \item The protocol exposes enough generic capability information for \IPP{} qualification, but no image-specific qualification rules.
\end{itemize}

\textbf{Gate.} A provider protocol is not accepted merely because it can allocate a buffer. It must have a coherent asynchronous lifecycle and a verifiable ownership/synchronization model.

\phase{3 --- ImageProcessingProtocol core implementation}

\textbf{Objective.} Implement the primary domain library against the already-defined generic execution contract.

\deliverables
\begin{enumerate}
  \item Concrete public types and descriptors for images, layers, pixel semantics, operation definitions, and typed operation inputs/outputs.
  \item Immutable operation instances with validation at construction or an explicit validated state.
  \item OperationImplementation protocol and registration metadata.
  \item Qualifying function types and structured qualification results that consume generic execution capabilities.
  \item Composition API for compatible operations.
  \item Explicit layer/composite targets and new-image/replacement result policy.
  \item Resource references that preserve the distinction between image semantics and generic execution allocation.
\end{enumerate}

\exitcriteria
\begin{itemize}
  \item The package imports \GEP{} in the permitted direction and the dependency graph is acyclic.
  \item Descriptor tests cover channel model, FP16/FP32/FP64 representation, strides, alpha, color metadata, non-color/scientific data, and invalid combinations.
  \item Operation tests prove immutable parameters and versioned semantic identity.
  \item Composition rejects incompatible input/output contracts.
  \item API tests prove that new-image derivation and replacement are distinct operations.
  \item No public API leaks provider-native resource classes.
\end{itemize}
